% RESEARCH_DRAFT_SHA256: 83db244b1ad84dff1730bc4ee8351b1cf3e8554767f3634811256634bab06e08
% Layout and source-reference wrappers expanded for the web renderer.
\documentclass[11pt]{article}
\usepackage[margin=25mm]{geometry}
\usepackage{amsthm}
\usepackage{fontspec,unicode-math}
\setmainfont{Latin Modern Roman}
\setmathfont{Latin Modern Math}
\usepackage{amsmath,amssymb,mathtools,xurl,hyperref}
\hypersetup{colorlinks=true,urlcolor=blue,linkcolor=blue}
\setcounter{secnumdepth}{0}
\setlength{\parindent}{0pt}
\setlength{\parskip}{5pt}
\setlength{\emergencystretch}{3em}
\newcommand{\sref}[2]{\hyperlink{source:#1:#2}{[S#2]}}
\newcommand{\cataloguescope}[1]{\paragraph{Source record.}#1}
\newcommand{\floor}[1]{\left\lfloor#1\right\rfloor}
\newtheorem{proposition}{Proposition}
\newtheorem{lemma}[proposition]{Lemma}
\begin{document}
\paragraph{Research manuscript.} OpenAI \texttt{gpt-6-astra}, reasoning \texttt{ultra}; 11 October 2026.
\paragraph{Preserved source body.} The complete catalogue mathematical text and sources follow with layout and source links expanded, including source qualifications and its historical preparation notes. New research begins at the section headed Research attempt.

% UnsolvedMath ID 20002121; source catalogue code AIM-INFRASTRUCTURE-0021
\section{A barrier-lattice audit for lecturer adoption of online assessment}\label{20002121}
{\small UnsolvedMath ID 20002121 · Computer Science}
\subsection{Definitions and mathematical statement}

\paragraph{Shared notation.}Standard number systems are $\mathbb N,\mathbb Z,\mathbb Q,\mathbb R,\mathbb C$. Specialized notation and index ranges follow the cited source and the definitions below.


Adoption is an instructor's decision to use an assessment tool; implementation is the way the tool is integrated into teaching. A barrier is a condition that inhibits adoption or effective use, such as unavailable support, incompatible workflows, or inaccessible content. Reported reasons and causal effects of removing a barrier are distinct kinds of evidence.
\paragraph{Record scope and qualifications.}
The catalogue title is retained as metadata. Named certificates, contracts, bounds or frameworks in that title are not asserted as results of the original source extract.
The extract does not identify a tool, institution, population, or operational measure of adoption. It does not posit a universal mathematical barrier model.

Catalogue screening limitation: Educational assessment, adoption, or study-design discussion; the supplied question does not state a mathematical target with defined assumptions and outcome.

The source snapshot labels this record partially solved; this is catalogue metadata rather than independent verification of a mathematical result or the current completion of a project.
\paragraph{Statement recorded in the supplied source.}
23. What stops other lecturers from using online assessment tools? \href{http://aimath.org/pastworkshops/oerassessmentproblems.pdf}{[S2]}
\paragraph{Source status.}
The supplied catalogue labels this source record partially solved; that metadata was checked against the supplied snapshot on 2026-10-11. The label is preserved for provenance and is not an independent verification of present open status or mathematical correctness.
\subsection{Short English statement}
Investigate barriers preventing lecturers from adopting online assessment tools.
\subsection{Sources}
\begin{itemize}
\item[S1] Ulam AI and the UnsolvedMath Contributors, UnsolvedMath: A Curated Collection of Open Mathematics Problems, record 20002121 (AIM-INFRASTRUCTURE-0021). Catalogue curation and metadata: CC BY 4.0. The source extract is retained; definitions, English summaries and qualifications are editorial additions.
\url{https://huggingface.co/datasets/ulamai/UnsolvedMath}.
\item[S2] AIM Workshop Problem Lists, source linked by catalogue record AIM-INFRASTRUCTURE-0021.
\url{http://aimath.org/pastworkshops/oerassessmentproblems.pdf}.
\end{itemize}
\paragraph{Preparation and review.}
The original source and UnsolvedMath contributors retain attribution for the source extract. Codex (GPT-6; exact serving revision not exposed) assisted with definitions, English summaries, scope notes and formatting. No new human mathematical authorship or independently verified solution is asserted. Permanent source ID 20002121 and code AIM-INFRASTRUCTURE-0021 checked against the complete supplied registry and current repository. Distinct upstream records, including repeated questions and broad or damaged extracts, are retained for complete coverage of the source CS category. A repeated record is not a new independent conjecture. Complete mathematical-equivalence and current-literature review has not been performed. Preparation and final review follow contribution prompts C and D, with source gaps disclosed. Structural and rendering checks do not establish a new theorem or open status.
\label{end:20002121}

\section{Research attempt}
\subsection{Literature, source context and scope}
The primary AIM notes describe a brainstorming session. The original excerpt and all catalogue qualifications remain above. They do not specify a mathematical theorem or a complete empirical study. The formal model below is an explicitly limited research route introduced for this attempt, not a reconstruction of a missing original conjecture. The supplied catalogue's earlier research summary was inspected as context; its named framework is not adopted as an established solution. The elementary result below is not claimed novel, and no educational effectiveness data were collected.

\subsection{Attempt: identifying barriers from support-bundle trials}
To examine the source's adoption-barrier question, I tested the information cost of learning an arbitrary deterministic monotone response to $n\ge2$ possible supports. A query chooses a support subset and observes only whether adoption occurs. Monotonicity says adding supports cannot turn adoption off. This is more general than a prespecified conjunction of barriers and is not an asserted instructor-behavior law.
\begin{proposition}
Any deterministic algorithm guaranteed to identify every such adoption function needs at least $\binom n{\lfloor n/2\rfloor}$ binary queries in the worst case, even with adaptively chosen bundles.
\end{proposition}
\begin{proof}
Put $k=\lfloor n/2\rfloor$. For every family $\mathcal S$ of $k$-element subsets, define $f_{\mathcal S}(A)=1$ exactly when $A$ contains some member of $\mathcal S$. Each function is monotone. On a $k$-element query set $A$, its value is one exactly when $A\in\mathcal S$, since containment between equal-size sets implies equality. Thus distinct families define distinct functions. There are $2^{\binom nk}$ such families.

An adaptive deterministic algorithm making at most $h$ binary queries has at most $2^h$ terminal transcripts. Two distinct functions cannot share a terminal transcript if the algorithm identifies both exactly: its output at that transcript is fixed. Hence $2^h\ge2^{\binom nk}$, giving $h\ge\binom nk$. The central binomial coefficient is at least $2^n/(n+1)$ because it is the largest of the $n+1$ nonnegative coefficients summing to $2^n$, so this bound is exponential up to a polynomial factor.
\end{proof}
\subsection{Verification and remaining gaps}
For six supports, the construction gives $2^{20}$ distinct functions and requires at least 20 binary queries in the worst case. Monotonicity alone therefore does not justify a small exact barrier-identification study. The theorem does not concern noisy population averages, approximate recovery, restricted behavioral families or which supports actually matter to lecturers. Those empirical and model-selection questions remain open in this attempt.

\section{References and source attribution}
American Institute of Mathematics, workshop discussion notes on open educational resources and assessment, \url{https://aimath.org/pastworkshops/oerassessmentproblems.pdf}. The 16-page primary PDF was opened on 11 October 2026 and its numbered brainstorming item was checked against the supplied record. The source supports the provenance and agenda, not the additional modeling assumptions or an empirical educational conclusion. Catalogue and contributor attribution are retained in the reproduced Sources section. The worked mathematics is proved here without reliance on an unverified external theorem.


\paragraph{Integration review.} Codex (GPT-6; exact serving revision not exposed) reviewed the written arguments and expanded layout and source-reference wrappers. Repository metadata and the production web renderer were checked using Python 3.13 and JavaScript with MathJax 3.2.2. This model-assisted review is not independent human proof verification; no PDF compilation or formal proof checking is asserted.

\section{Completion Estimate}
\noindent\textbf{COMPLETION: 0\%}
No substantiated new progress toward the entire original target is claimed. The worked results and failed route are delimited above; a known result, illustrative model, or auxiliary obstruction is not counted as original-target completion.

\section{Final status and attribution}
\textbf{Attempt status: unresolved.} This describes this research attempt, not a new ruling on the source catalogue's status or on external literature claims.
The mathematics and this self-audit were generated by OpenAI \texttt{gpt-6-astra}, reasoning \texttt{ultra}, on 11 October 2026. Codex (GPT-6; exact serving revision not exposed) orchestrated the batch and contributed formatting/repository review. No human mathematical authorship, independent human proof checking, or formal verification is claimed. The historical preparation paragraphs in the reproduced catalogue body concern that body only.

Contribution prompt A was applied to the complete definition and supplied source record, followed by the adversarial checks of prompt B. Local reproduction files are \texttt{batch008/checks.py}, \texttt{computations.json}, and the per-ID audit JSON. They were not public links at preparation time. Numerical checks are finite and do not prove unrestricted results. No LaTeX engine was found locally; compilation was not performed by this research agent.

\end{document}
